%% =====================================================================
%%  B1-T-06-01 -- what B1 contributes to "Protest Too Much"
%%  -------------------------------------------------------------------
%%  LAYER A: APPEND-ONLY. Written once, then left alone. Material found
%%  only here sits in the `unique` slot and is protected by rule R10.
%% =====================================================================
%% @kind: source
%% @id: B1-T-06-01
%% @book: B1
%% @topic: T-06-01
%% @locator: Tactic 1, pp. 9--18
%% @status: distilled
%% @unique: yes
%% @integrated: yes
%% @updated: 2026-09-19

\dossier{B1}{T-06-01}{\bookshortB1}{Tactic 1, pp. 9--18}

\begin{distilled}
  \item People who stress something too much are doubted by the experienced observer; the pattern is drawn from Shakespeare's protest-too-much observation and matched against the author's own losses.
  \item A liar soothe-talks in a too-reassuring, too-demonstrative way, over-compensating nervously for a lack of substance.
  \item The author lists the exact reassurances used by the people who defrauded him, each delivered immediately before the loss.
  \item Terms that genuinely favour the other party are said not to require selling; the proposer adopts a take-it-or-leave-it stance instead.
  \item The same dynamic is said to apply outside money: a person who must remind you of their trustworthiness is probably using you.
\end{distilled}

\begin{unique}
  \item The over-compensation mechanism --- nervous excess of emphasis substituting for absent substance --- stated as the reason the tell works, rather than as a folk observation.
  \item Extension of the tell from financial dealings to personal relationships, where the same self-advertising of reliability appears.
\end{unique}

\begin{terms}
  \item \emph{protest too much} --- asserting a quality so insistently that the assertion becomes evidence against it.
\end{terms}
